\section{Introduction}

\subsection*{Motivation}

Priority queues arise naturally wherever a common server must satisfy heterogeneous demand. The primary motivation for this thesis is healthcare operations, specifically the management of waiting lists for elderly-care facilities, where medical severity dictates a user's —hereafter \emph{customer}— priority class. Here prolonged waiting may itself alter the system's dynamics: a lack of appropriate care can deteriorate a waiting customer's health, leading them to change priority status (jockeying) or leave the system (abandonment), whether to seek alternative care or through death. Despite this motivation, the work is fundamentally modelling-driven: we characterise the long-run behaviour of queueing systems in which jockeying and/or abandonment play a role.

This thesis analyses a $2$-class $M/M/1$ queue with non-preemptive priority: each class arrives according to a Poisson process (rates $\lambda_1$ and $\lambda_2$), service times are exponential at a common rate $\mu$, and the single server completes the current service even when a higher-priority customer arrives. To this we add the mechanisms of jockeying and abandonment. The model and its mechanisms are described in more detail in \S\ref{sec:model_x}.

How these two mechanisms interact within a priority discipline, and how that interaction affects the difficulty of obtaining a closed-form Probability Generating Function (PGF), is the main focus of this thesis.

\subsection*{The analytical challenge}

In full generality, the bivariate PGF for Model-$X$ satisfies a functional equation with terms of three kinds: algebraic coefficients, partial derivatives in $x$ —from jockeying and abandonment in class-$1$— and partial derivatives in $y$ —from the same mechanisms in class-$2$. When both kinds of derivative term are present, the equation belongs to a family of integro-differential equations whose solution would require more advanced tools. The problem is moreover asymmetric: derivative terms from the non-priority class complicate it more than those from the priority class alone.

We therefore solve less complex models obtained by setting some of the jockeying and abandonment parameters to zero. Table~\ref{tab:model_overview_intro} lists the models studied here and indicates whether a PGF expression is obtained for each:
\begin{table}[H]
\centering
\renewcommand{\arraystretch}{1.3}
\label{tab:model_overview_intro}
    \begin{tabular}{clll}
        \hline
            \textbf{Model} & \textbf{Active mechanism} & \textbf{Equation type} & \textbf{Status} \\
            \hline
            A     & none                                    & algebraic PGF equation & solved  \\
            B     & bidirectional jockeying                 & PDE, Volterra family   & open \\
            X     & jockeying and abandonment               & PDE, transcendental characteristics & open \\
            B$_2$   & one-way jockeying ($1\to2$)           & ODE in $x$, Kummer     & solved \\
            $\BH$ & one-way jockeying, head-of-line       & algebraic PGF equation & solved \\
            C$_2$   & class-1 abandonment                   & ODE in $x$, Kummer     & solved \\
            $\CH$ & class-1 abandonment, head-of-line     & algebraic PGF equation & solved \\
        \hline
    \end{tabular}
    \caption{Models studied in this thesis, with the active mechanism, the resulting functional-equation type, and the solution status. The head-of-line variants $\BH$ and $\CH$ restrict the mechanism to the customer at the head of Queue~1, which
collapses the derivative term to an algebraic one and renders the model solvable in closed
form.}
\end{table}
Some models —notably Model-$X$ and Model-$B$— are left incomplete because they involve derivative terms in $y$; others, such as a potential Model-$B_1$, Model-$C$ or Model-$C_1$, are not considered for the same reason, as they all reduce to a Volterra integral equation lying beyond the scope of this thesis. Model-$B$ serves as our representative example of how such problems reduce to that integral equation; the remaining models are solvable. The last two entries, Models~$\BH$ and~$\CH$ (\S\S\ref{sec:model_b21} and~\ref{sec:model_c21}), are \emph{head-of-line} simplifications in which only the customer at the head of Queue~1 may jockey or abandon. The mechanism then acts at the constant rate $\gamma_1\mathbf 1_{\{n_1\ge1\}}$ or $\theta_1\mathbf 1_{\{n_1\ge1\}}$ rather than the length-proportional $\gamma_1 n_1$ or $\theta_1 n_1$, so the fundamental equation stays algebraic and is solved by the Kernel Method exactly as in Model-$A$.

\subsection*{Methods}

Three complementary proof strategies are used throughout.

\emph{Kernel method / ODE with integrating factor}. For Model-$A$, which has no derivative terms, the algebraic coefficient of $P(x,y)$ vanishes on a curve (the \emph{kernel}); evaluating the equation there determines the remaining unknowns. Model-$B_2$ and Model-$C_2$, which do carry derivative terms in $x$, reduce to a first-order linear ODE in $x$ with an explicit integrating factor, and their unknowns follow from a regularity argument at a singularity of that factor. We present both strategies together, as they share the same underlying idea: the latter extends the former to the case where the unknowns are independent of the derivative terms and can be taken outside the integral.

\emph{Probabilistic argument via Pollaczek--Khinchine and the effective service time}. For Model-$A$ and Model-$C_2$, class-$2$ arrivals are Poisson, so the PASTA property holds and we can apply the Pollaczek--Khinchine formula together with the \emph{effective service time} seen by a class-$2$ customer. This is an $M/G/1$ queue whose service times equal the busy period of the class-$1$ subsystem —an $M/M/1$ queue for Model-$A$ and an $M/M/1+M$ system for Model-$C_2$. In both cases the argument independently confirms the analytic results. Since jockeying destroys the Poisson structure, it does not apply to the other models.

\emph{Partial PGF (PPGF) in the alternative state space $\widetilde{S}$}. For the baseline Model-$A$ we introduce an alternative state space indexed by the number of class-$2$ customers —whose distribution is hard— and the total number of customers —which is easy, since the total is conserved and behaves like a regular $M/M/1$ queue. This yields a PPGF with a non-homogeneous term. Cohen's recursion argument \cite{cohen1956delay} was developed for the regular state space; we explore, with limited success, the challenges of carrying it over to this alternative space and assess whether Cohen's trick can deliver a feasible rough approximation.

\subsection*{Structure of the thesis}

Section~\ref{sec:literature} reviews the relevant literature and Section~\ref{sec:prelim} collects the probabilistic and analytic background used throughout this work. Section~\ref{sec:model_x} defines the full model, Model-$X$, and derives the general balance equations and the fundamental PGF equations, both for the regular and the alternative state space. The analysis then proceeds from baseline to synthesis: Section~\ref{sec:model_a} solves the baseline Model-$A$ in closed form; Section~\ref{sec:model_b} shows that bidirectional jockeying (Model-$B$) reduces the problem to a Volterra-type integral equation and is left open. Sections~\ref{sec:model_b2} and~\ref{sec:model_c2} then recover tractability in the one-way-jockeying (Model-$B_2$) and class-1-abandonment (Model-$C_2$) specialisations, while Sections~\ref{sec:model_b21} and~\ref{sec:model_c21} show how the head-of-line variants, Models~$\BH$ and~$\CH$, collapse the functional equation to an algebraic one solved in closed form. Finally, Section~\ref{sec:comparison} compares the solved models, Section~\ref{sec:results} presents the numerical validation and performance metrics, and Section~\ref{sec:conclusion} concludes by abstracting the obstruction taxonomy and setting out the future work.
